\documentclass[12pt]{article}
\usepackage[T1]{fontenc}
\usepackage[utf8]{inputenc}
\usepackage{lmodern}
\usepackage{textcomp}
\usepackage{enumitem}
\usepackage{geometry}
\geometry{margin=1in}
\begin{document}
\begin{center}
Answer Key - Geography - K-12 Level Assessment 18 \
\small (Answers are ordered exactly as in the question paper.)
\end{center}

\section*{Multiple Choice}
\begin{enumerate}[label=\arabic*.]
\item \textbf{Answer:} A
\\ \textit{Options:} A) the addition of non-basic workers to an urban economy that has added more basic workers.; B) a ratio of urban to rural workers in an economy.; C) the ratio of city workers to farm workers in an economy.; D) the addition of rural jobs created by the urban economy.
\\ \textit{Note:} The correct answer describes the multiplier effect as it highlights how an increase in basic workers, who provide essential goods and services, leads to a proportional increase in non-basic workers that support and sustain the urban economy, thereby amplifying economic growth.
\item \textbf{Answer:} B
\\ \textit{Options:} A) CBR - CDR 5 2.; B) Births 1 immigration 5 deaths 1 emigration.; C) Total fertility 5 3.; D) Replacement fertility is achieved.
\\ \textit{Note:} Zero population growth (ZPG) occurs when the number of births plus immigration equals the number of deaths plus emigration, resulting in no net change in the population size.
\item \textbf{Answer:} D
\\ \textit{Options:} A) Longitude; B) Latitude; C) Equator; D) Distance to the nearest city
\\ \textit{Note:} Distance to the nearest city is not a method for determining absolute location because absolute location is defined by specific coordinates (latitude and longitude), rather than relative distances to nearby landmarks.
\item \textbf{Answer:} D
\\ \textit{Options:} A) Phosphate; B) Oil; C) Petroleum; D) Coal
\\ \textit{Note:} Coal is considered the most abundant fossil fuel because it is found in greater quantities worldwide compared to oil and natural gas, making it a primary energy source in many countries.
\item \textbf{Answer:} C
\\ \textit{Options:} A) global grid system.; B) Prime Meridian; C) latitude.; D) longitude.
\\ \textit{Note:} Latitude is the geographic coordinate that measures the distance north or south of the equator, which is defined as 0 degrees latitude.
\end{enumerate}

\section*{True / False}
\begin{enumerate}[label=\arabic*.]
\setcounter{enumi}{5}
\item \textbf{Answer:} True
\\ \textit{Options:} A) True; B) False
\\ \textit{Note:} The statement is true because continental margins provide access to vital resources such as fish, minerals, and fertile land, as well as trade routes, which attract large populations seeking economic opportunities and sustainable livelihoods.
\item \textbf{Answer:} False
\\ \textit{Options:} A) True; B) False
\\ \textit{Note:} The statement is false because while better schools can attract families, they are not typically classified as pull factors in migration studies, which generally focus on economic, social, and environmental conditions rather than educational quality alone.
\item \textbf{Answer:} True
\\ \textit{Options:} A) True; B) False
\\ \textit{Note:} An urban area encompasses both the densely populated central city and the adjacent suburbs that are socially and economically linked to it, thus justifying the statement as true.
\item \textbf{Answer:} False
\\ \textit{Options:} A) True; B) False
\\ \textit{Note:} A 'hearth' refers to the origin or source area where a cultural trait begins, not the destination point where it ends up after diffusion.
\item \textbf{Answer:} True
\\ \textit{Options:} A) True; B) False
\\ \textit{Note:} Islam is the most rapidly growing religion in the United States due to factors such as high birth rates among Muslim families, immigration, and increasing conversion rates.
\end{enumerate}

\section*{Long Form Response}
\begin{enumerate}[label=\arabic*.]
\setcounter{enumi}{10}
\item \textbf{Answer:} The Untouchables, also known as Dalits, play a significant role in the context of the Hindu caste system, which categorizes people into hierarchical groups. Traditionally, Untouchables are considered outside the four main castes and have faced severe social exclusion and discrimination. This exclusion impacts social structure by creating deep divisions within society, limiting access to education, employment, and basic rights for millions of people. Geographically, Untouchables often reside in separate areas, reinforcing their marginalization and highlighting disparities in development across different regions of India. Their situation illustrates the ongoing challenges of caste- based discrimination and its effects on both social cohesion and geographic distribution in the country.
\\ \textit{Note:} The answer correctly identifies the Untouchables as a marginalized group within the Hindu caste system, emphasizing their social exclusion and its profound effects on societal divisions, opportunities, and geographic segregation in India.
\item \textbf{Answer:} Transnational companies (TNCs) are large corporations that operate in multiple countries, often with a headquarters in one nation and branches or subsidiaries in others. They are characterized by their ability to leverage resources and markets globally, which allows them to maximize profits and minimize costs. The belief that TNCs are controlled by foreign governments is inaccurate because TNCs are primarily driven by their corporate interests and profit motives rather than external political influences. While they may engage with governments in various countries, their operational decisions are made independently to serve their shareholders and business strategies, not dictated by foreign governments. This independence underscores the complex relationship TNCs have with national economies, emphasizing their role as global entities rather than extensions of governmental control.
\\ \textit{Note:} correct because it emphasizes that TNCs operate independently based on their profit motives and corporate interests, rather than being influenced or controlled by foreign governments.
\end{enumerate}

\vfill
\noindent\textit{End of Answer Key}
\end{document}